\section{Indices e Otimizacao}

\subsection{Estrategia de Indexacao}

\subsubsection{Indices Primarios}

\begin{lstlisting}[caption=Indices essenciais para performance]
-- Indices para consultas frequentes por usuario
CREATE INDEX idx_transacoes_usuario_data_status 
    ON transacoes(usuario_id, data_transacao DESC, status);

CREATE INDEX idx_contas_usuario_ativa_tipo 
    ON contas(usuario_id, ativa, tipo);

-- Indices para consultas de grupo
CREATE INDEX idx_transacoes_grupo_data 
    ON transacoes(grupo_id, data_transacao DESC) 
    WHERE grupo_id IS NOT NULL;

-- Indices para relatorios
CREATE INDEX idx_transacoes_categoria_data 
    ON transacoes(categoria_id, data_transacao DESC, status);

-- Indices para busca textual
CREATE INDEX idx_transacoes_descricao_gin 
    ON transacoes USING gin(to_tsvector('portuguese', descricao));
\end{lstlisting}

\subsubsection{Indices Condicionais}

\begin{lstlisting}[caption=Indices otimizados com condicoes]
-- Apenas transacoes ativas
CREATE INDEX idx_transacoes_ativas 
    ON transacoes(usuario_id, data_transacao DESC) 
    WHERE status = 'confirmada';

-- Apenas contas ativas
CREATE INDEX idx_contas_ativas_saldo 
    ON contas(usuario_id, saldo_atual) 
    WHERE ativa = true;

-- Apenas usuarios verificados
CREATE INDEX idx_usuarios_verificados 
    ON usuarios(email) 
    WHERE email_verificado = true AND ativo = true;

-- Recorrencias ativas
CREATE INDEX idx_recorrencias_proxima 
    ON recorrencias(proxima_execucao) 
    WHERE ativa = true;
\end{lstlisting}

\subsection{Monitoramento de Performance}

\subsubsection{Views de Monitoramento}

\begin{lstlisting}[caption=Views para analise de performance]
-- View para consultas lentas
CREATE VIEW slow_queries AS
SELECT 
    query,
    calls,
    total_time,
    mean_time,
    rows,
    100.0 * shared_blks_hit / nullif(shared_blks_hit + shared_blks_read, 0) AS hit_percent
FROM pg_stat_statements
WHERE calls > 100
ORDER BY mean_time DESC;

-- View para uso de indices
CREATE VIEW index_usage AS
SELECT 
    schemaname,
    tablename,
    indexname,
    idx_tup_read,
    idx_tup_fetch,
    idx_scan,
    pg_size_pretty(pg_relation_size(indexrelid)) as size
FROM pg_stat_user_indexes
ORDER BY idx_scan DESC;

-- View para tabelas grandes
CREATE VIEW table_sizes AS
SELECT 
    schemaname,
    tablename,
    pg_size_pretty(pg_total_relation_size(schemaname||'.'||tablename)) as size,
    pg_total_relation_size(schemaname||'.'||tablename) as bytes
FROM pg_tables
WHERE schemaname = 'public'
ORDER BY bytes DESC;
\end{lstlisting}

\subsection{Otimizacoes Avancadas}

\subsubsection{Particao de Tabelas}

\begin{lstlisting}[caption=Estrategia de particao para transacoes]
-- Particionar transacoes por ano para melhor performance
CREATE TABLE transacoes_2024 PARTITION OF transacoes 
    FOR VALUES FROM ('2024-01-01') TO ('2025-01-01');

CREATE TABLE transacoes_2025 PARTITION OF transacoes 
    FOR VALUES FROM ('2025-01-01') TO ('2026-01-01');

-- Funcao para criar particoes automaticamente
CREATE OR REPLACE FUNCTION criar_particao_mensal()
RETURNS void AS $$
DECLARE
    ano INTEGER := EXTRACT(YEAR FROM CURRENT_DATE);
    mes INTEGER := EXTRACT(MONTH FROM CURRENT_DATE);
    data_inicio DATE;
    data_fim DATE;
    nome_tabela TEXT;
BEGIN
    data_inicio := DATE_TRUNC('month', CURRENT_DATE);
    data_fim := data_inicio + INTERVAL '1 month';
    nome_tabela := 'transacoes_' || ano || '_' || LPAD(mes::TEXT, 2, '0');
    
    EXECUTE format('CREATE TABLE IF NOT EXISTS %I PARTITION OF transacoes 
                    FOR VALUES FROM (%L) TO (%L)',
                   nome_tabela, data_inicio, data_fim);
END;
$$ LANGUAGE plpgsql;
\end{lstlisting}

\subsubsection{Cache e Buffers}

\begin{lstlisting}[caption=Configuracoes recomendadas para PostgreSQL]
-- Configuracoes no postgresql.conf
-- shared_buffers = 256MB (25% da RAM)
-- effective_cache_size = 1GB (75% da RAM)
-- work_mem = 4MB
-- maintenance_work_mem = 64MB
-- checkpoint_completion_target = 0.9
-- wal_buffers = 16MB
-- default_statistics_target = 100

-- Query para verificar configuracoes atuais
SELECT name, setting, unit, context 
FROM pg_settings 
WHERE name IN (
    'shared_buffers',
    'effective_cache_size', 
    'work_mem',
    'maintenance_work_mem'
);
\end{lstlisting}

\subsection{Analise e Manutencao}

\subsubsection{Scripts de Analise}

\begin{lstlisting}[caption=Scripts para analise de performance]
-- Atualizar estatisticas
ANALYZE;

-- Verificar bloat em tabelas
SELECT 
    tablename,
    pg_size_pretty(pg_total_relation_size(tablename::regclass)) as size,
    round(100 * (pg_total_relation_size(tablename::regclass) - 
                 pg_relation_size(tablename::regclass))::numeric / 
                 pg_total_relation_size(tablename::regclass), 2) as bloat_pct
FROM pg_tables 
WHERE schemaname = 'public'
  AND pg_total_relation_size(tablename::regclass) > 1000000
ORDER BY pg_total_relation_size(tablename::regclass) DESC;

-- Verificar indices nao utilizados
SELECT 
    schemaname,
    tablename,
    indexname,
    idx_scan,
    pg_size_pretty(pg_relation_size(indexname::regclass)) as size
FROM pg_stat_user_indexes
WHERE idx_scan = 0
  AND schemaname = 'public'
ORDER BY pg_relation_size(indexname::regclass) DESC;
\end{lstlisting}

\subsubsection{Manutencao Automatica}

\begin{lstlisting}[caption=Funcao para manutencao automatica]
CREATE OR REPLACE FUNCTION manutencao_automatica()
RETURNS void AS $$
BEGIN
    -- Vacuum analyze em tabelas principais
    VACUUM ANALYZE usuarios;
    VACUUM ANALYZE grupos;
    VACUUM ANALYZE contas;
    VACUUM ANALYZE categorias;
    VACUUM ANALYZE transacoes;
    
    -- Refresh views materializadas
    REFRESH MATERIALIZED VIEW CONCURRENTLY resumo_financeiro_mensal;
    
    -- Limpar logs antigos (>90 dias)
    DELETE FROM audit.audit_logs 
    WHERE timestamp_operacao < CURRENT_DATE - INTERVAL '90 days';
    
    -- Limpar sessoes expiradas
    DELETE FROM sessoes 
    WHERE data_expiracao < CURRENT_TIMESTAMP;
    
    -- Log da manutencao
    INSERT INTO audit.audit_logs (tabela, operacao, dados_novos)
    VALUES ('sistema', 'MAINTENANCE', 
            jsonb_build_object('timestamp', CURRENT_TIMESTAMP, 'status', 'concluida'));
END;
$$ LANGUAGE plpgsql;
\end{lstlisting}
